\scanpage{22}
\mspage{9}
\[
(10)\quad\left\{\begin{aligned}
\rho^3&=1\\
\varepsilon_\infty^2\varepsilon_1^2\varepsilon_0^2&=1\\
\deleted{\unsure{\varepsilon_0}}\ \varepsilon_1&=\varepsilon_0^{\deleted{4}}\varepsilon_\infty^4\rho\varepsilon_0^2\deleted{\rho}\varepsilon_\infty^2\\
\varepsilon_\infty&=\varepsilon_1\varepsilon_0^4\rho\varepsilon_1^2\deleted{\rho}\varepsilon_0^2\\
\varepsilon_0&=\varepsilon_\infty\varepsilon_1^4\rho\varepsilon_\infty^2\deleted{\rho}\varepsilon_1^2.
\end{aligned}\right.
\]
\margin{il faut aussi ajouter (au moins!) les relations $\rho\varepsilon_0\rho^{-1}=\varepsilon_1$ etc donc on a deux gén. $(\varepsilon_0,\rho)$.}
On peut encore se demander si $\rho$ s'exprime en fonction des $\varepsilon_i$, et on trouve comme tantôt que non.

Quant à relever $\dot\tau$ en un élément de $\mathcal E$, par un chemin de $P$ à $\bar P$, on ne voit pas de choix privilégié --- on voit trois choix également naturels. Au lieu de les expliciter d'emblée, je préfère considérer le sous-groupe d'isotropie de $P$ dans $\Gamma=\Gamma_{0,3}=\unsure{\mathfrak S_3|\mathcal S_3}\times\mathbb Z/2$, qui est formé des couples $\dot\sigma,\dot\tau^\varepsilon$ tels que $\operatorname{sg}(\dot\sigma)=\varepsilon$; \deleted{On} i.e. les éléments
\[
\widetilde{\dot\sigma}=(\dot\sigma,\dot\tau^{\operatorname{sg}(\dot\sigma)})\in\Gamma
\]
qui sont relevés canoniquement en des éléments $(\widetilde{\dot\sigma},\text{lacet trivial en }P)$ de $\mathcal E$. On trouve un hom.
\[
(11)\quad\left\{\begin{aligned}
\varphi:\unsure{\mathfrak S_3|\mathcal S_3}&\longrightarrow\Gamma=\unsure{\mathfrak S_3|\mathcal S_3}\times\mathbb Z/2\\
\dot\sigma&\longmapsto\dot\sigma,\operatorname{sg}(\dot\sigma)
\end{aligned}\right.
\]
et un relèvement canonique de $\varphi$ en
\[
(12)\qquad\sigma\longrightarrow\tilde\sigma:\quad\unsure{\mathfrak S_3|\mathcal S_3}\longrightarrow\mathcal E.
\]
On pose
\[
\tilde\sigma_i=\widetilde{\dot\sigma_i}
\]
\margin{on a d'autre part \deleted{$\unsure{\tilde\rho}$} $\widetilde{\dot\rho}=\rho$}
\mspage{21}
